% arara: clean: {extensions: ['log','aux','synctex.gz','out']}
% arara: pdflatex: { shell: yes }
% arara: pdflatex: { shell: yes }
% arara: clean: {extensions: ['log','aux','synctex.gz','out']}
\documentclass[spanish,addpoints,answers,a4paper]{exam}
\usepackage[T1]{fontenc}
\usepackage[spanish]{babel}
\usepackage{libertine}
\usepackage[tmargin=2cm,rmargin=2cm,lmargin=2cm]{geometry}
\usepackage{amsmath,bm}
\usepackage[shortlabels]{enumitem}
\usepackage{siunitx}
\sisetup{
	group-four-digits = true,
	group-separator = {,}
}
\usepackage{multicol}

\usepackage{graphicx}
\graphicspath{{../../../images/}}
\pagestyle{plain}

\usepackage{minted}
\setminted[cpp]{tabsize=2}
\usepackage{hyperref}
\renewcommand{\solutiontitle}{\noindent\textbf{Solución}\par\noindent}
\renewcommand\listingscaption{Listado}
\newcommand{\unmarkedfntext}[1]{%
	\begingroup
	\renewcommand\thefootnote{}\footnote{#1}%
	\addtocounter{footnote}{-1}%
	\endgroup
}
\spanishdatedel
\decimalpoint

\def\LOGOUNI{%
	\begin{picture}(0,0)\unitlength=1cm
	\put (-13.5,-2.5) {\includegraphics[width=2.8cm]{logouni}}
	\end{picture}
}

\def\LOGOCTIC{%
	\begin{picture}(0,0)\unitlength=1cm
	\put (-8,-1.2) {\includegraphics[height=1cm]{logocticblack}}
	\end{picture}
}

\newcommand{\mychar}{%
	\begingroup\normalfont
	\includegraphics[height=\fontcharht\font`\B]{Octocat.png}%
	\endgroup
}

\begin{document}

\begin{center}
	\sffamily\bfseries\scshape
	{\Large UNIVERSIDAD NACIONAL DE INGENIERÍA}\LOGOUNI\\
	Centro de Tecnologías de la Información y Comunicaciones\LOGOCTIC\\
\end{center}

\vspace{.8cm}

\begin{center}\sffamily\bfseries\large
	Programación en \texttt{C++} \\
	TEMA $\bm{1}$: Entrada y salida de datos
\end{center}

\vspace{.5cm}
\noindent
\makebox[\textwidth]{Nombres y apellidos:\enspace\hrulefill}
\vspace{0.2in}
\makebox[\textwidth]{Nombres y apellidos del instructor:\enspace MSc. Víctor Melchor Espinoza.\hfill}

\section*{Ejercicios de ejemplo}

\subsection*{Primera parte}

\begin{questions}

\question Convertir en expresiones aritméticas algorítmicas las siguientes expresiones algebraicas:
\begin{multicols}{2}
	\begin{parts}
			\part $5(x+y)$
			\part $a^2+b^2$
			\part $\frac{x+y}{u+\frac{v}{w}}$
			\part $\frac{x}{y}(v+w)$
	\end{parts}
\end{multicols}

\begin{solution}

\end{solution}

\question Escribir las siguientes expresiones algorítmicas como expresiones algebraicas:
\begin{multicols}{2}
	\begin{parts}
			\part \texttt{b**2 - 4*a*c}
			\part \texttt{3*x**4 - 5*x**3 + 12*x - 17}
			\part \texttt{(b + d)/(c + 4)}
			\part \texttt{(x**2 + y**2)**(1/2)}
	\end{parts}
\end{multicols}

\begin{solution}

\end{solution}

\question Determine los resultados que se obtienen al aplicar la función correspondiente:
\begin{multicols}{2}
	\begin{parts}
			\part \texttt{raizcua(25)}
			\part \texttt{redondeo(6.5)}
			\part \texttt{redondeo(3.1)}
			\part \texttt{redondeo(-3.2)}
			\part \texttt{trunc(5.6)}
			\part \texttt{trunc(3.1)}
			\part \texttt{trunc(-3.8)}
			\part \texttt{cuadrado(7)}
			\part \texttt{abs(19.3)}
			\part \texttt{abs(-15)}
	\end{parts}
\end{multicols}

\begin{solution}

\end{solution}

\question ¿Cuáles son los valores de \mintinline{cpp}{A}, \mintinline{cpp}{B} y \mintinline{cpp}{C} después de las siguientes operaciones?

\begin{minted}{cpp}
A = 3;
B = 4;
C = A + 2*B;
B = C - A;
A = B*C;
\end{minted}

\begin{solution}

\end{solution}

\question ¿Cuál es el valor de \mintinline{cpp}{x} después de las siguientes operaciones?

\begin{minted}{cpp}
x = 2;
x = cuadrado(x + x);
x = raizcua(16 + raizcua(x) + 5);
\end{minted}

\begin{solution}

\end{solution}

\question Encontrar el valor de la variable \mintinline{cpp}{RESUL} después de la ejecución de las siguientes operaciones:
\begin{multicols}{2}
\begin{parts}
\part
\begin{minted}{cpp}
X = 3;
Y = 2;
RESUL = X**Y - Y
\end{minted}
\part
\begin{minted}{cpp}
RESUL = 5;
X = 3;
RESUL = RESUL*X + RESUL - X;
\end{minted}
\end{parts}
\end{multicols}

\begin{solution}

\end{solution}

\question ¿Cuál es el resultado de evaluar las siguientes expresiones?
\begin{multicols}{3}
\begin{parts}
	\part \mintinline{cpp}{3 + 6*14 - 18/3;}
	\part \mintinline{cpp}{8 + 7*3 - 4*6;}
	\part \mintinline{cpp}{-4*7 + 2**3/4 - 5;}
\end{parts}
\end{multicols}
\question Calcule el valor de las siguientes expresiones:
\begin{multicols}{2}
\begin{parts}
	\part \mintinline{cpp}{8 + 7*3 + 4*5;}
	\part \mintinline{cpp}{-2**3 + 16 DIV 5;}
	\part \mintinline{cpp}{(32 + 3*4)/5 - (19 MOD 7) DIV 3;}
	\part \mintinline{cpp}{(2**2*3)- 3 MOD 5 + 4^3;}
	\part \mintinline{cpp}{3 + 2*(18 - 4**2);}
	\part \mintinline{cpp}{16*6 - 3*2 -12;}
\end{parts}
\end{multicols}

\begin{solution}

\end{solution}

\end{questions}

\subsection*{Segunda parte}

Para los siguientes ejercicios, elabore los algoritmos en diagrama de flujo y luego en pseudocódigo. Recuerde que debe usar la sintáxis y el estilo establecido en el curso.

\begin{questions}

\question Determinar el área de un cuadrado conociendo su lado.
\question Determinar el área de un triángulo conociendo su base y su altura.
\question Determinar el área de un trapecio conociendo sus bases y su altura.
\question Convertir grados \textbf{Centígrados} a grados \textbf{Farenheit} a partir de la siguiente fórmula:

{\centering \mintinline{cpp}{GradosFar = 1.8*GradosCent + 32}\par}

\begin{solution}

\end{solution}

\question Una señora de una tienda compra un artículo a determinado precio, y lo vende con una ganancia del 18\%. Determinar el precio final de venta.

\begin{solution}

\end{solution}

\question Todos los lunes, miércoles y viernes; una persona corre la misma ruta y cronometra los tiempos obtenidos. Determinar el tiempo promedio que la persona tarda en recorrer la ruta en una semana cualesquiera.

\begin{solution}

\end{solution}

\question Calcule e imprima el número de segundos que hay en número \mintinline{cpp}{X} de días.

\begin{solution}

\end{solution}

\question En una determinada gasolinera, los surtidores miden y registran en galones, pero el precio de la gasolina está en litros. Calcular e imprimir lo que hay que cobrarle al cliente en soles.

\textbf{Consideraciones}: Cada galón tiene \SI{3.785}{\litre} y \SI[per-mode=symbol]{15}[S/]{\per\litre}.

\end{questions}

\section*{Ejercicios de práctica}

\begin{questions}

\question Elabore un programa que envíe a pantalla el saludo ``\textbf{Hola mundo}''.

\begin{solution}
	\begin{listing}[H]
		\footnotesize
		\inputminted{cpp}{1.cc}
		\caption{Programa \texttt{1.cc}.}
		\label{lst:1.1}
	\end{listing}
\end{solution}

\question Presente un ejemplo del uso de los operadores de asignación y aritméticos.

\begin{solution}
	\begin{listing}[H]
		\footnotesize
		\inputminted{cpp}{2.cc}
		\caption{Programa \texttt{2.cc}.}
		\label{lst:1.2}
	\end{listing}
\end{solution}

\question Determinar el salario neto de un trabajador que trabajó $36$ horas, con un pago por hora de \SI{50.00}[S/]{} soles y la tasa del impuesto del $13$\%.\label{q:3}

\begin{solution}
	\begin{listing}[H]
		\footnotesize
		\inputminted{cpp}{3.cc}
		\caption{Programa \texttt{3.cc}.}
		\label{lst:1.3}
	\end{listing}
\end{solution}

\question Resolver el ejercicio~\ref{q:3} para el caso general, conociendo el número de horas trabajadas, pago por hora y la tasa de impuesto.

\begin{solution}
	\begin{listing}[H]
		\footnotesize
		\inputminted{cpp}{4.cc}
		\caption{Programa \texttt{4.cc}.}
		\label{lst:1.4}
	\end{listing}
\end{solution}

\question Escriba un programa que pida dos números enteros que representen al dividendo y divisor. Obtenga el cociente y el resto.

\begin{solution}
	\begin{listing}[H]
		\footnotesize
		\inputminted{cpp}{5.cc}
		\caption{Programa \texttt{5.cc}.}
		\label{lst:1.5}
	\end{listing}
\end{solution}

\question Dado un número de tres cifras, hallar la suma de sus cifras.

\begin{solution}
	\begin{listing}[H]
		\footnotesize
		\inputminted{cpp}{6.cc}
		\caption{Programa \texttt{6.cc}.}
		\label{lst:1.6}
	\end{listing}
\end{solution}

\question Pida un valor del radio y obtenga el área del círculo.

\begin{solution}
		\begin{listing}[H]
		\footnotesize
		\inputminted{cpp}{7.cc}
		\caption{Programa \texttt{7.cc}.}
		\label{lst:1.7}
	\end{listing}
\end{solution}

\question Elabore un programa que pida un monto en dólares y obtenga el monto equivalente en soles. Defina el tipo de cambio como una constante.

\begin{solution}
	\begin{listing}[H]
		\footnotesize
		\inputminted{cpp}{8.cc}
		\caption{Programa \texttt{8.cc}.}
		\label{lst:1.8}
	\end{listing}
\end{solution}

\end{questions}

\section*{Ejercicios obligatorios}

\begin{questions}

\question En el curso de \textbf{Programación}, el promedio final se compone de los siguientes porcentajes:

\begin{itemize}[--]
	\item $60$\% del promedio de cuatro exámenes (PE).
	\item $15$\% del promedio de laboratorios (PL)
	\item $15$\% del promedio de tareas (PT)
	\item $10$\% de la nota del proyecto final (PR)
\end{itemize}

Desarrolle un programa en el que solicite las notas de los cuatro exámenes, del promedio de laboratorios, del promedio de tareas y la nota del proyecto. En base a esto, el programa determinará la nota final del alumno.

\textbf{Ejemplo de salida:}
\begin{minted}{bash}
Bienvenido al programa que calcula la Nota Final del curso de Programación.
Ingrese la nota del Examen N°1: 10
Ingrese la nota del Examen N°2: 08
Ingrese la nota del Examen N°3: 09
Ingrese la nota del Examen N°4: 09
Ingrese el promedio de Laboratorios: 10
Ingrese el promedio de Tareas: 14
Ingrese la nota del Proyecto Final: 14

Reporte de notas:
Nota del Examen N°1: 10
Nota del Examen N°2: 08
Nota del Examen N°3: 09
Nota del Examen N°4: 09
Promedio de Exámenes (PE): 09
Promedio de Laboratorios (PL): 14
Promedio de Tareas (PT): 14
Nota del Proyecto Final (PF): 14
La fórmula de la Nota Final es: NF=0.60PE+0.15P.L+0.15PT+0.10PR
Vuestra Nota Final es: 11
\end{minted}

\begin{solution}
	\begin{listing}[H]
		\footnotesize
		\inputminted{cpp}{1.cpp}
		\caption{Programa \texttt{1.cpp}.}
		\label{lst:1.9}
\end{listing}
\end{solution}

\question Realiza un programa que proporcione el desglose en billetes y monedas de una cantidad entera de euros. Recuerda que hay billetes de $500$, $200$, $100$, $50$, $20$, $10$ y $5$ euros, también hay monedas de $2$ y $1$ euro.

\begin{minted}{bash}
Ingrese el importe total (en euros) = 977

1 billete de 500 euros
2 billetes de 200 euros
1 billete de 50 euros
1 billete de 20 euros
0 billete de 10 euros
1 moneda de 5 euros
2 monedas de 1 euro
\end{minted}

\begin{solution}
	\begin{listing}[H]
		\footnotesize
		\inputminted{cpp}{2.cpp}
		\caption{Programa \texttt{2.cpp}.}
		\label{lst:1.10}
\end{listing}
\end{solution}

\end{questions}

\begin{flushright}\bfseries
Centro de Tecnologías de la Información y Comunicaciones (CTIC)\\[2mm]
\today\unmarkedfntext{Código disponible en \href{https://github.com/carlosal1015/Cpp-Programming}{\mychar{}}.}
\end{flushright}

\end{document}